\documentclass[../../main.tex]{subfiles}
\begin{document}

{\bf [解]}
$P,Q$が線分$OR$上にある時を時刻の基準として$t=0$とする．
又この時の$R$の座標が$(3,0)$となるようにして考える．
すると題意より時刻$t$での各々の座標は
\begin{align}
P(\cos t,\sin t),\quad Q(2\cos2t,2\sin2t),\quad R(3\cos3t,3\sin3t)
\end{align}
とかける．以下$s=\sin t,\ c=\cos t$とする．


\begin{figure}[htb]
\begin{center}
\begin{tikzpicture}[scale=1.1, >=stealth]
  % 1. 座標軸
  \draw[->] (-2.3,0) -- (3.5,0) node[right] {$x$};
  \draw[->] (0,-2.3) -- (0,3.5) node[above] {$y$};

  % 2. 半径 1,2,3 の同心円 (P,Q,Rの軌道)
  \draw (0,0) circle (1);
  \draw (0,0) circle (2);
  \draw (0,0) circle (3);
  \node[below right] at (1,0) {$1$};
  \node[below right] at (2,0) {$2$};
  \node[below right] at (3,0) {$3$};
  \fill (0,0) circle (1.1pt) node[below left] {$O$};

  % 3. P,Q,R の座標定義
  \coordinate (P) at (0.878,0.479);
  \coordinate (Q) at (1.081,1.683);
  \coordinate (R) at (0.212,2.992);

  % 4. △PQR の斜線塗りつぶしと外枠描画
  % pattern=north east lines で右上上がりの斜線（角度や密度も調整可能）
  \fill[pattern=north east lines, pattern color=black!60] (P) -- (Q) -- (R) -- cycle;
  \draw[thick] (P) -- (Q) -- (R) -- cycle;

  % 5. 点のプロットとラベル
  \fill (P) circle (1.1pt) node[below right] {$P$};
  \fill (Q) circle (1.1pt) node[right] {$Q$};
  \fill (R) circle (1.1pt) node[above] {$R$};
  
\end{tikzpicture}
\end{center}
\caption{点$P$，$Q$，$R$の位置関係}
\end{figure}


この時
\begin{align}
\overrightarrow{PQ}=\begin{pmatrix}2\cos2t-c\\2\sin2t-s\end{pmatrix},\qquad
\overrightarrow{PR}=\begin{pmatrix}3\cos3t-c\\3\sin3t-s\end{pmatrix}
\end{align}
だからサラスの公式より，時刻$t$での$\triangle PQR$の面積$T(t)$として，
\begin{align}
\begin{split}
2T(t)&=\left|(3\sin3t-s)(2\cos2t-c)-(3\cos3t-c)(2\sin2t-s)\right| \\
&=\bigl|6\sin3t\cos2t-2s\cos2t-3c\sin3t+sc-(6\sin2t\cos3t-3s\cos3t-2c\sin2t+sc)\bigr| \\
&=\bigl|6(\sin3t\cos2t-\sin2t\cos3t)-3(c\sin3t-s\cos3t)-2(s\cos2t-c\sin2t)\bigr| \\
&=\bigl|6\sin t-3\sin2t-2\sin(-t)\bigr|
=\bigl|8s-3\sin2t\bigr|
=\bigl|8s-6sc\bigr| \quad \cdots \text{①}
\end{split} \\
\therefore 
T(t) &= \bigl|4s-3sc\bigr| \label{eq:1}
\end{align}
と書ける．\cref{eq:1}の右辺の絶対値の中身を$f(t)$としてその挙動を求めれば良い．
一階微分は
\begin{align}
f'(t)=8c-6\cos2t=8c-6(2c^2-1)=2(-6c^2+4c+3)
\end{align}
$f'(t)=0$となるのは
\begin{align}
 t_{\pm} = \dfrac{2\pm \sqrt{22}}{6}
\end{align}
だから増減表は以下のようになる．ここで$|s|,c$の周期性から，$0\le t\le\pi$で考えれば良い．

\begin{table}[htb]
\begin{center}
\caption{$f(t)$の増減表}
\begin{tabular}{c|ccccc}
$t$ & $0$ & & & & $\pi$ \\
\hline
$c$ & $1$ & & $t_{-}$ & & $-1$ \\
\hline
$f'$ & & $+$ & $0$ & $-$ & \\
\hline
$f$ & $0$ & $\nearrow$ & & $\searrow$ & $0$
\end{tabular}
\end{center} 
\end{table}

従って求める$T(t)$の最大値は$t=t_{-}$の時である．この時$\sin t$は
\begin{align}
|s|=\sqrt{1-t_{-}^2}=\dfrac{1}{6}\sqrt{10+4\sqrt{22}}
\end{align}
だから，\cref{eq:1}に代入して求める最大値は
\begin{align}
\max T(t) 
 &= T(t_{-}) \\
 &=\dfrac{\sqrt{10+4\sqrt{22}}}{6}\left|4-3\cdot\dfrac{2-\sqrt{22}}{6}\right| \\
 &=\dfrac{(6+\sqrt{22})\sqrt{10+4\sqrt{22}}}{12} \\
 &=\dfrac{\sqrt{409+88\sqrt{22}}}{6}
\end{align}
である．


\newpage

\end{document}
